%======================================================================
\section{Electronegative Ambipolar Diffusion Correction}
\label{sec:electronegative-ambipolar}
%======================================================================

The prior Phase-1 release of this framework carried an architectural
limitation (L1, see \S\ref{sec:hybrid-limits} of the prior revision) in
which the 2D ambipolar diffusion solver used the electropositive
coefficient
\[
  D_a^\mathrm{ep} \;=\; D_i\,\Bigl(1 + \tfrac{T_e}{T_i}\Bigr),
\]
valid only in the limit of negligible negative-ion density.  The present
revision resolves L1 by replacing this expression with its
electronegative generalisation due to Lieberman \& Lichtenberg
(2005)~\cite{Lieberman2005} Ch.~10.3:
\begin{equation}
  \boxed{\;D_a^\mathrm{en} \;=\; D_i\,\Bigl(1 + \tfrac{T_e}{T_i}\Bigr)
  \;\frac{1 + \alpha}{1 + \alpha\,T_i/T_e}\;,\;}
  \label{eq:Da-en-L1}
\end{equation}
where $\alpha = n_{-}/n_e$ is the electronegativity.  The correction
factor $(1+\alpha)/(1+\alpha T_i/T_e)$ reduces exactly to unity when
$\alpha \to 0$, so the electropositive form is recovered in the appropriate
limit.

\subsection{Rigorous Physics, Not a Calibration}

This modification is not a calibration.  Unlike the two fitted parameters
$\gamma_\mathrm{Al} = 0.18$ (wall-recombination probability) and
$\lambda_\mathrm{exp} = 3.20$ (bias-sheath expansion factor), the
electronegative correction introduces \emph{no} free parameters.  The
three inputs $\alpha$, $T_e$, and $T_i$ are all obtained from independent
first-principles solves already present in the framework:
\begin{itemize}
  \item $\alpha$ from the 0D global model's attachment--recombination
  balance (\S\ref{sec:0d-2d-coupling}),
  \item $T_e(r,z)$ from the local power-balance solve in m11,
  \item $T_i = k_B T_\mathrm{gas}$ with $T_\mathrm{gas} = 313$\,K taken
  as the wall-set operating temperature (not tuned).
\end{itemize}
Equation~\ref{eq:Da-en-L1} is the algebraically correct transport
coefficient for a plasma with three charged species (electrons, a
dominant positive ion, and a dominant negative ion); the prior
electropositive form dropped the last two terms and was numerically
convenient but physically incomplete.

A dedicated unit test
(\texttt{tests/test\_electronegative\_ambipolar.py}) verifies the three
invariants that characterise the correction factor:
\begin{itemize}
  \item $f(\alpha{=}0) = 1$ exactly — the electropositive limit is bit-for-bit recovered.
  \item $f(T_i{=}T_e) = 1$ exactly — the isothermal degenerate case.
  \item $f(\alpha{=}1,\,T_e{=}20\,\text{eV},\,T_i{=}1\,\text{eV})
  = 2/1.05 \approx 1.905$ — the Lieberman worked example.
\end{itemize}
All seven tests pass (together with the eight self-consistent-$\eta$ tests
inherited from the prior revision).

\subsection{Implementation}

The change is surgical: a single block in
\texttt{src/dtpm/solvers/ambipolar\_diffusion.py} now constructs
$D_a^\mathrm{en}$ via Eq.~\ref{eq:Da-en-L1}.  The scalar
$\alpha$ from the 0D closure is read once per Picard iteration from
\texttt{result\_0D['alpha']} (file \texttt{m11\_plasma\_chemistry.py}) and
is threaded uniformly into the 2D solver.  A spatially resolved
$\alpha(r,z)$ field (Phase-3 scope) would be a refinement; the scalar
form is exact for the spatially averaged 0D closure that is this
framework's hybrid architecture.

The generic diffusion-matrix assembler
(\texttt{build\_diffusion\_matrix()} in
\texttt{src/dtpm/solvers/species\_transport.py}) was \emph{not} modified:
it accepts any $D_\mathrm{field}$ as input, so the new $D_a^\mathrm{en}$
propagates through the existing finite-difference stencil without change.
The module-level \texttt{multispecies\_transport.py} neutral-species
transport is independent of $\alpha$ and was also not touched.

\subsection{Quantitative Impact at the Mettler Operating Point}
\label{sec:L1-finding-alpha-small}

An important physics finding emerges from applying Eq.~\ref{eq:Da-en-L1}
at Mettler's canonical operating point (\SI{1000}{\watt}, \SI{10}{\milli\torr},
90\,\% SF$_6$, \SI{200}{\watt} bias).  The 0D global model returns
\[
  \alpha_{0\mathrm{D}} \;\approx\; 0.02,
\]
which is \emph{two orders of magnitude} below the typical-electronegative
value ($\alpha \sim 1$--$1.5$) cited as a generic textbook estimate in the
prior revision's L1 discussion.  The cause is the very high electron density
at Mettler's high-power / high-bias operating point
($n_e \sim 10^{19}$\,m$^{-3}$), which replenishes the electron population
faster than the dissociative-attachment rate depletes it: $\alpha$ is
determined by the dynamical balance $R_\mathrm{att}/k_\mathrm{rec}/n_e$
(\texttt{global\_model.py} Eqs.~202--204), and at these conditions the
denominator $k_\mathrm{rec} n_e$ dominates.

The correction factor evaluated at $\alpha = 0.02$ and $T_i/T_e \approx
0.012$ is
\[
  \frac{1+\alpha}{1+\alpha\,T_i/T_e} \;\approx\; 1.019\,,
\]
so the electronegative $D_a^\mathrm{en}$ is only $\sim 2\,\%$ above the
electropositive $D_a^\mathrm{ep}$ in this regime.  The predicted
centre-to-edge [F]-drop changes by ${<}1$\,percentage point between the
two solvers.

The prior revision's projection that L1 would close the approximately
6-percentage-point gap between the modelled 68\,\% and the experimental
74\,\% [F] drop was therefore based on an over-estimate of $\alpha$.  That
projection is now corrected.  The residual gap between the model and
Mettler's absolute data at this operating point must be attributed to
effects other than L1 --- most likely an over-tuned
$\gamma_\mathrm{Al} = 0.18$ (fitted against Mettler's 74\,\% on the
approximate D$_a^\mathrm{ep}$ solver) and/or the composition-insensitive
$\gamma_\mathrm{Al}$ architectural choice itself (see remaining
limitations in \S\ref{sec:hybrid-limits}).

\subsection{Regime Where L1 Matters}

The correction factor $(1+\alpha)/(1+\alpha T_i/T_e)$ approaches 2 only
when $\alpha T_i/T_e$ is much smaller than 1 \emph{and} $\alpha \gtrsim 1$.
At higher pressure (lower $n_e$ at fixed power) or lower ICP power (lower
$n_e$ at fixed pressure), $\alpha$ rises and the correction becomes
substantial.  Example: at \SI{50}{\milli\torr}/\SI{300}{\watt} (outside the
present Mettler validation window), the 0D solver returns $\alpha \approx
1.4$, giving a correction factor of $\sim 2.3$.  The L1 machinery is
therefore physically important for higher-pressure / lower-power regimes
and remains a required component of the framework even though its
quantitative effect at Mettler's 1000 W point is minor.  This is a
significant methodological finding: \emph{applying a rigorous physics
correction does not always produce a large numerical change}, and it is
only by implementing the correction that one can bound its quantitative
relevance at a specific operating point.
